\documentclass[aspectratio=169, 9pt]{beamer}
\usetheme{metropolis}

\usepackage{xeCJK}
\setCJKmainfont{PingFang SC}
\setCJKsansfont{PingFang SC}
\setCJKmonofont{PingFang SC}

\usepackage{fontspec}
\setmonofont{Menlo}[Scale=0.82]
\usepackage{tikz}
\usetikzlibrary{arrows.meta, positioning, calc, shapes.geometric, shadows.blur}
\usepackage{graphicx}
\usepackage{booktabs}
\usepackage{array}
\usepackage{listings}
\usepackage{alltt}
\usepackage{xcolor}

\graphicspath{{images/}}

% ── Color palette ──────────────────────────────────────────────────
\definecolor{primary}{HTML}{6C5CE7}
\definecolor{accent}{HTML}{00B894}
\definecolor{warning}{HTML}{E17055}
\definecolor{dark}{HTML}{2D3436}
\definecolor{light}{HTML}{DFE6E9}
\definecolor{muted}{HTML}{636E72}
\definecolor{codebg}{HTML}{F7F7FA}
\definecolor{erlred}{HTML}{D63031}
\definecolor{erlblue}{HTML}{0984E3}
\definecolor{nodegreen}{HTML}{00A884}
\definecolor{terminalfg}{HTML}{102A3A}
\definecolor{terminalbg}{HTML}{FBFBFE}
\definecolor{terminalborder}{HTML}{DFE6E9}

% ── Beamer theme ───────────────────────────────────────────────────
\setbeamercolor{normal text}{fg=dark,bg=white}
\setbeamercolor{frametitle}{fg=white,bg=primary}
\setbeamercolor{progress bar}{fg=accent,bg=light}
\setbeamercolor{alerted text}{fg=warning}
\setbeamercolor{block title}{fg=white,bg=primary}
\setbeamercolor{block body}{fg=dark,bg=primary!7}
\setbeamercolor{block title example}{fg=white,bg=accent}
\setbeamercolor{block body example}{fg=dark,bg=accent!7}
\setbeamercolor{block title alerted}{fg=white,bg=warning}
\setbeamercolor{block body alerted}{fg=dark,bg=warning!8}
\setbeamertemplate{itemize item}{\color{accent}$\blacktriangleright$}
\setbeamertemplate{itemize subitem}{\color{primary}$\bullet$}
\setbeamertemplate{enumerate item}{\color{accent}\insertenumlabel.}
\metroset{progressbar=frametitle, numbering=fraction}

% ── Macros ─────────────────────────────────────────────────────────
\newcommand{\keypoint}[1]{\textbf{\color{primary}#1}}
\newcommand{\good}[1]{\textbf{\color{accent!80!black}#1}}
\newcommand{\warn}[1]{\textbf{\color{warning}#1}}
\newcommand{\bad}[1]{\textbf{\color{erlred}#1}}
\newcommand{\code}[1]{\texttt{#1}}
\newcommand{\imgslide}[3]{%
  \begin{frame}{#1}
    \begin{columns}[T,onlytextwidth]
      \column{0.60\textwidth}
      \centering
      \includegraphics[width=\linewidth,height=0.72\textheight,keepaspectratio]{#2}
      \column{0.36\textwidth}
      #3
    \end{columns}
  \end{frame}
}
\newcommand{\wideimgslide}[3]{%
  \begin{frame}{#1}
    \centering
    \includegraphics[width=\linewidth,height=0.66\textheight,keepaspectratio]{#2}
    \vspace{0.15cm}
    #3
  \end{frame}
}

\lstdefinelanguage{Erlang2}{
  keywords={module,export,application,supervisor,gen_server,gen_statem,case,of,end,fun,receive,after,true,false,ok},
  sensitive=true,
  morecomment=[l]{\%},
  morestring=[b]",
  morestring=[b]'
}
\lstset{
  language=Erlang2,
  basicstyle=\ttfamily\scriptsize,
  keywordstyle=\color{primary}\bfseries,
  commentstyle=\color{muted}\itshape,
  stringstyle=\color{accent!70!black},
  backgroundcolor=\color{codebg},
  frame=single,
  framesep=1.5mm,
  rulecolor=\color{light},
  showstringspaces=false,
  breaklines=true,
  tabsize=2
}

\title{Application（一）runtime application}
\subtitle{OTP 系统的启动边界}
\author{}
\date{}

\begin{document}

% ── Title ──────────────────────────────────────────────────────────
\begin{frame}[plain]
  \begin{tikzpicture}[remember picture, overlay]
    \fill[primary!8] (current page.north west) rectangle (current page.south east);
    \fill[primary!15] (current page.north west) ++(1.7,-1.3) circle (2.7cm);
    \fill[accent!12]  (current page.south east) ++(-2.1,1.6) circle (3.0cm);
    \node[anchor=center, font=\fontsize{25}{31}\selectfont\bfseries, text=dark]
      at ([yshift=1.75cm]current page.center) {Application（一）runtime application};
    \node[anchor=center, font=\Large, text=primary!80]
      at ([yshift=0.55cm]current page.center) {OTP 系统的启动边界};
    \node[anchor=center, font=\large\bfseries, text=accent!80!black]
      at ([yshift=-0.35cm]current page.center) {从 rebar3 配置，到 Observer 里的 supervision tree};
    \draw[accent, line width=1.5pt]
      ([yshift=-1.15cm, xshift=-4.8cm]current page.center)
      -- ([yshift=-1.15cm, xshift=4.8cm]current page.center);
    \node[anchor=south, font=\scriptsize, text=dark!45]
      at ([yshift=0.4cm]current page.south)
      {理解路径：工程结构 \enspace|\enspace .app.src \enspace|\enspace application:start/stop \enspace|\enspace Observer 验证};
  \end{tikzpicture}
\end{frame}

% ── Concepts ───────────────────────────────────────────────────────
\begin{frame}{Application 先分成两种}
  \begin{columns}[T,onlytextwidth]
    \column{0.48\textwidth}
    \begin{block}{runtime application}
      \begin{itemize}
        \item 也叫 \keypoint{active application}
        \item 也叫 \keypoint{process-owning application}
        \item 启动后创建自己的 process
        \item 通常启动 top supervisor
        \item Observer 中能看到 supervision tree
      \end{itemize}
    \end{block}
    \column{0.48\textwidth}
    \begin{exampleblock}{library application}
      \begin{itemize}
        \item 也叫 \keypoint{passive application}
        \item 也叫 \keypoint{code-only application}
        \item 主要提供 module 和 function
        \item 通常没有自己的长期进程
        \item 下一集用 math 作为例子
      \end{itemize}
    \end{exampleblock}
  \end{columns}
  \vspace{0.25cm}
  \centering
  \good{两者都是 OTP application，区别在于是否拥有自己的运行时进程结构。}
\end{frame}

\begin{frame}[fragile]{runtime application 的心智模型}
  \begin{columns}[T,onlytextwidth]
    \column{0.48\textwidth}
    \begin{block}{从配置到运行时}
\begin{lstlisting}
.app / .app.src
  -> mod 字段
      -> application callback
          -> start/2
              -> top supervisor
                  -> worker process
\end{lstlisting}
    \end{block}
    \column{0.48\textwidth}
    \begin{alertblock}{这一集的重点}
      \begin{itemize}
        \item application 的实现很复杂
        \item 但用户入口非常集中
        \item 大多数时候只改配置和启动命令
        \item rebar3 已经包装了编译、依赖、启动
        \item Observer 展示启动后的真实结果
      \end{itemize}
    \end{alertblock}
  \end{columns}
  \vspace{0.2cm}
  \centering
  \warn{我们讲原理，是为了看懂现象；不是要求一开始背 application controller 的内部实现。}
\end{frame}

% ── Screenshot sequence ────────────────────────────────────────────
\begin{frame}[fragile]{第一张图：工程本身已经是 OTP 结构}
  \begin{columns}[T,onlytextwidth]
    \column{0.54\textwidth}
    \begin{block}{截图文字化：\code{tree}}
\begin{lstlisting}[
  basicstyle=\ttfamily\tiny\color{terminalfg},
  backgroundcolor=\color{terminalbg},
  frame=single,
  rulecolor=\color{terminalborder},
  framesep=2.2mm,
  keepspaces=true,
  columns=fullflexible,
  escapeinside={(*@}{@*)}
]
[ericksun@ERICKSUN-MC1:full_otp_app] (main) $ tree
.
|-- README.txt
|-- (*@\textcolor{erlblue}{data}@*)
|-- rebar.config
`-- (*@\textcolor{erlblue}{src}@*)
    |-- data_cache.erl
    |-- data_store.erl
    |-- data_store_sup.erl
    |-- demo.erl
    |-- event_bus.erl
    |-- event_bus_handler.erl
    |-- full_otp_app.app.src
    |-- full_otp_app_app.erl
    |-- full_otp_app_sup.erl
    |-- kv_server.erl
    `-- traffic_light.erl

3 directories, 13 files
\end{lstlisting}
    \end{block}

    \column{0.42\textwidth}
    \begin{block}{截图解析}
      \begin{itemize}
        \item 根目录有 \code{rebar.config}，说明项目由 rebar3 管理。
        \item \code{src/} 下有 \code{full\_otp\_app.app.src}、\code{full\_otp\_app\_app.erl}、\code{full\_otp\_app\_sup.erl}。
        \item 这三个文件正好对应 runtime application 的三层入口：resource file、callback module、top supervisor。
      \end{itemize}
    \end{block}
  \end{columns}
\end{frame}

\begin{frame}[fragile]{第二张图：rebar3 把启动入口包装在配置里}
  \begin{columns}[T,onlytextwidth]
    \column{0.50\textwidth}
    \begin{block}{rebar.config}
\begin{lstlisting}[
  basicstyle=\ttfamily\tiny\color{terminalfg},
  backgroundcolor=\color{terminalbg},
  frame=single,
  rulecolor=\color{terminalborder},
  framesep=2.2mm,
  keepspaces=true,
  columns=fullflexible,
  escapeinside={(*@}{@*)},
  language={}
]
[ericksun@ERICKSUN-MC1:full_otp_app] (main) $ (*@\textcolor{accent}{cat}@*) rebar.config
{(*@\textcolor{erlblue}{erl\_opts}@*), [debug_info]}.
{(*@\textcolor{erlblue}{deps}@*), [
    {(*@\textcolor{accent}{observer\_cli}@*), "1.8.8"},
    {(*@\textcolor{accent}{recon}@*), "2.5.6"}
]}.
{(*@\textcolor{erlblue}{shell}@*), [
    {(*@\textcolor{accent}{apps}@*), [(*@\textcolor{accent}{full\_otp\_app}@*)]}
]}.
\end{lstlisting}
    \end{block}
    \column{0.46\textwidth}
    \begin{block}{依赖}
      \code{observer\_cli} 和 \code{recon} 被声明为 dependencies，构建时会被拉取和编译。
    \end{block}
    \vspace{0.3cm}
    \begin{exampleblock}{shell 配置}
      \code{\{apps, [full\_otp\_app]\}} 表示进入 shell 时自动启动这个 application。
    \end{exampleblock}
  \end{columns}
\end{frame}

\begin{frame}[fragile]{第三张图：compile 阶段已经在分析 application}
  \begin{columns}[T,onlytextwidth]
    \column{0.56\textwidth}
    \begin{block}{rebar3 compile 输出}
\begin{lstlisting}[
  basicstyle=\ttfamily\tiny\color{terminalfg},
  backgroundcolor=\color{terminalbg},
  frame=single,
  rulecolor=\color{terminalborder},
  framesep=2.2mm,
  keepspaces=true,
  columns=fullflexible,
  escapeinside={(*@}{@*)},
  language={}
]
[ericksun@ERICKSUN-MC1:full_otp_app] (main) $ (*@\textcolor{accent}{rebar3 compile}@*)
(*@\textcolor{erlblue}{Verifying dependencies...}@*)
(*@\textcolor{warning}{Fetching observer\_cli v1.8.8}@*)
(*@\textcolor{warning}{Fetching recon v2.5.6}@*)
(*@\textcolor{erlblue}{Analyzing applications...}@*)
(*@\textcolor{warning}{Compiling recon}@*)
(*@\textcolor{warning}{Compiling observer\_cli}@*)
===> (*@\textcolor{erlblue}{Analyzing applications...}@*)
(*@\textcolor{warning}{Compiling full\_otp\_app}@*)
\end{lstlisting}
    \end{block}
    \column{0.40\textwidth}
    \begin{block}{截图解析}
      \begin{itemize}
        \item rebar3 先 verifying dependencies，再 fetching 依赖。
        \item 之后出现 \code{Analyzing applications...}。
        \item 最后编译 \code{full\_otp\_app} 本身。
      \end{itemize}
    \end{block}
  \end{columns}
  \vspace{0.1cm}
  \centering
  \good{这说明 application 不是临时概念，而是构建工具认识的工程单元。}
\end{frame}

\begin{frame}[fragile]{第四张图：真正入口在 .app.src 的 mod 字段}
  \begin{columns}[T,onlytextwidth]
    \column{0.54\textwidth}
    \begin{block}{src/full\_otp\_app.app.src}
\begin{lstlisting}[
  basicstyle=\ttfamily\tiny\color{terminalfg},
  backgroundcolor=\color{terminalbg},
  frame=single,
  rulecolor=\color{terminalborder},
  framesep=1.8mm,
  keepspaces=true,
  columns=fullflexible,
  escapeinside={(*@}{@*)},
  language={}
]
[ericksun@ERICKSUN-MC1:full_otp_app] (main) $ (*@\textcolor{accent}{cat}@*) src/full_otp_app.app.src
{(*@\textcolor{erlblue}{application}@*), full_otp_app,
 [{description, "...gen_server, gen_event and gen_statem"},
  {(*@\textcolor{primary}{vsn}@*), "0.1.0"},
  {(*@\textcolor{erlblue}{registered}@*), [
      full_otp_app_sup,
      kv_server,
      traffic_light,
      event_bus,
      data_store_sup,
      data_store,
      data_cache
  ]},
  {(*@\textcolor{warning}{mod}@*), {full_otp_app_app, []}},
  {(*@\textcolor{erlblue}{applications}@*), [
      kernel,
      stdlib,
      observer_cli,
      recon
  ]},
  {env, []},
  {modules, []},
  {licenses, ["Apache-2.0"]},
  {links, []}
 ]}.
\end{lstlisting}
    \end{block}
    \column{0.42\textwidth}
    \begin{block}{resource file 里最关键的三类信息}
      \begin{itemize}
        \item \code{registered}: 这个 app 注册的名字
        \item \code{mod}: 启动 callback module
        \item \code{applications}: 依赖哪些 application
      \end{itemize}
    \end{block}
    \vspace{0.15cm}
    \begin{alertblock}{runtime application 的关键线索}
      \centering
      {\ttfamily \{mod, \{full\_otp\_app\_app, []\}\}}
      \par\vspace{0.15cm}
      有 \code{mod}，就意味着 application controller 启动它时，会调用 callback module。
    \end{alertblock}
  \end{columns}
\end{frame}

\begin{frame}[fragile]{callback module 通常只做一件事}
  \begin{columns}[T,onlytextwidth]
    \column{0.52\textwidth}
    \begin{block}{application callback}
\begin{lstlisting}
-module(full_otp_app_app).
-behaviour(application).

-export([start/2, stop/1]).

start(_StartType, _StartArgs) ->
    full_otp_app_sup:start_link().

stop(_State) ->
    ok.
\end{lstlisting}
    \end{block}
    \column{0.44\textwidth}
    \begin{exampleblock}{分工非常清晰}
      \begin{itemize}
        \item application callback 不写业务逻辑
        \item 它只把系统交给 top supervisor
        \item 之后长期运行的是 supervisor 和 worker
        \item 所以 application 是启动边界，不是 main 函数
      \end{itemize}
    \end{exampleblock}
  \end{columns}
\end{frame}

\begin{frame}[fragile]{第五张图：shell 启动后，full\_otp\_app 已经在列表里}
  \begin{columns}[T,onlytextwidth]
    \column{0.60\textwidth}
    \begin{block}{rebar3 shell + application:which\_applications().}
\begin{lstlisting}[
  basicstyle=\ttfamily\fontsize{5}{6}\selectfont\color{terminalfg},
  backgroundcolor=\color{terminalbg},
  frame=single,
  rulecolor=\color{terminalborder},
  framesep=1.8mm,
  keepspaces=true,
  columns=fullflexible,
  escapeinside={(*@}{@*)},
  language={}
]
[ericksun@ERICKSUN-MC1:full_otp_app] (main) $ (*@\textcolor{accent}{rebar3 shell --sname app\_demo}@*)
(*@\textcolor{erlblue}{Verifying dependencies...}@*)
(*@\textcolor{erlblue}{Analyzing applications...}@*)
(*@\textcolor{warning}{Compiling full\_otp\_app}@*)
Erlang/OTP 26 [erts-14.2.5.12] [64-bit] [smp:12:12] [jit]

Eshell V14.2.5.12 (press Ctrl+G to abort, type help(). for help)
(*@\textcolor{accent}{[full\_otp\_app] starting...}@*)
(*@\textcolor{accent}{[kv\_server] init: pid=<0.258.0>}@*)
(*@\textcolor{accent}{[event\_bus\_handler] init}@*)
(*@\textcolor{accent}{[event\_bus] started with default handler}@*)
(*@\textcolor{accent}{[traffic\_light] init: starting at red}@*)
(*@\textcolor{accent}{[data\_store] init: pid=<0.262.0>}@*)
(*@\textcolor{accent}{[data\_cache] init: ETS cache table created}@*)
===> (*@\textcolor{accent}{Booted full\_otp\_app}@*)
(app_demo@ERICKSUN-MC1)1> (*@\textcolor{primary}{application:which\_applications().}@*)
[{(*@\textcolor{erlblue}{full\_otp\_app}@*), "...gen_server, gen_event and gen_statem", "0.1.0"},
 {observer_cli, "Visualize Erlang Nodes On The Command Line", "1.8.8"},
 {recon, "Diagnostic tools for production use", "2.5.6"},
 {inets, "INETS  CXC 138 49", "9.1.0.4"},
 {ssl, "Erlang/OTP SSL application", "11.1.4.10"},
 {crypto, "CRYPTO", "5.4.2.3"},
 {stdlib, "ERTS  CXC 138 10", "5.2.3.5"},
 {(*@\textcolor{primary}{kernel}@*), "ERTS  CXC 138 10", "9.2.4.10"}]
\end{lstlisting}
    \end{block}
    \column{0.36\textwidth}
    \begin{block}{截图解析}
      \begin{itemize}
        \item \code{rebar3 shell --sname} 启动节点。
        \item 日志显示 \code{kv\_server}、\code{event\_bus}、\code{traffic\_light}、\code{data\_store}、\code{data\_cache} 被初始化。
        \item \code{application:which\_applications()} 中出现 \code{full\_otp\_app}。
      \end{itemize}
    \end{block}
    \vspace{0.15cm}
    \centering
    \good{从 shell 角度看，application 已经启动；从进程角度看，一批 worker 已经存在。}
  \end{columns}
\end{frame}

\imgslide{第六张图：Observer 把启动结果画成 supervision tree}{06-observer-app.png}{
  \begin{block}{左侧 application 列表}
    \code{full\_otp\_app}、\code{inets}、\code{kernel}、\code{ssl} 都是当前节点中的 application。
  \end{block}
  \begin{exampleblock}{右侧运行时结构}
    \code{full\_otp\_app\_sup} 下挂着 \code{kv\_server}、\code{event\_bus}、\code{traffic\_light} 和 \code{data\_store\_sup}。
  \end{exampleblock}
}

\begin{frame}{把前六张图串起来}
  \begin{block}{runtime application 的启动链路}
    \centering
    \renewcommand{\arraystretch}{1.45}
    \begin{tabular}{@{}>{\centering\arraybackslash}p{0.16\linewidth}>{\centering\arraybackslash}p{0.04\linewidth}>{\centering\arraybackslash}p{0.18\linewidth}>{\centering\arraybackslash}p{0.04\linewidth}>{\centering\arraybackslash}p{0.18\linewidth}>{\centering\arraybackslash}p{0.04\linewidth}>{\centering\arraybackslash}p{0.18\linewidth}@{}}
      \keypoint{工程结构} & $\rightarrow$ & \keypoint{工具配置} & $\rightarrow$ & \keypoint{.app.src} & $\rightarrow$ & \keypoint{callback} \\
      {\ttfamily src/} & & {\ttfamily rebar.config} & & {\ttfamily mod} & & {\ttfamily start/2} \\
    \end{tabular}
    \vspace{0.2cm}
    \begin{tabular}{@{}>{\centering\arraybackslash}p{0.32\linewidth}>{\centering\arraybackslash}p{0.08\linewidth}>{\centering\arraybackslash}p{0.32\linewidth}@{}}
      \keypoint{top supervisor} & $\rightarrow$ & \keypoint{Observer} \\
      {\ttfamily full\_otp\_app\_sup} & & {\ttfamily supervision tree} \\
    \end{tabular}
  \end{block}
  \begin{alertblock}{这就是 runtime application 的核心闭环}
    \centering
    \Large 配置文件给出入口，callback 启动 supervisor，Observer 展示最终运行结构。
  \end{alertblock}
\end{frame}

\begin{frame}[fragile]{第七张图：application 可以被 stop 管理}
  \begin{columns}[T,onlytextwidth]
    \column{0.56\textwidth}
    \begin{block}{shell 中执行 stop}
\begin{lstlisting}[
  basicstyle=\ttfamily\scriptsize\color{terminalfg},
  backgroundcolor=\color{terminalbg},
  frame=single,
  rulecolor=\color{terminalborder},
  framesep=2.2mm,
  keepspaces=true,
  columns=fullflexible,
  escapeinside={(*@}{@*)},
  language={}
]
(app_demo@ERICKSUN-MC1)3> (*@\textcolor{primary}{application:stop(full\_otp\_app).}@*)
(*@\textcolor{accent}{[event\_bus\_handler] terminate: reason=stop}@*)
(*@\textcolor{muted}{[full\_otp\_app] stopped.}@*)
=INFO REPORT==== ... ===
    application: full_otp_app
    exited: stopped
    type: temporary
(*@\textcolor{accent}{ok}@*)
\end{lstlisting}
    \end{block}
    \column{0.40\textwidth}
    \begin{block}{截图解析}
      \begin{itemize}
        \item shell 中执行 \code{application:stop(full\_otp\_app).}
        \item 业务日志打印 \code{[full\_otp\_app] stopped.}
        \item INFO REPORT 显示 application \code{full\_otp\_app} exited: stopped。
        \item 这说明 application 是可启动、可停止、可管理的 OTP 单元。
      \end{itemize}
    \end{block}
  \end{columns}
\end{frame}

\imgslide{第八张图：stop 之后，Observer 中 app 消失}{08-observer-app-stop.png}{
  \begin{alertblock}{Observer 证明了 stop 的效果}
    左侧 application 列表里已经没有 \code{full\_otp\_app}，右侧也不再有 \code{full\_otp\_app\_sup} 这棵 supervision tree。
  \end{alertblock}
  \vspace{0.2cm}
  \good{application:stop/1 管理的不是一个函数调用，而是一整棵运行时结构的生命周期。}
}

\begin{frame}[fragile]{第九张图：application:start 重新建立进程树}
  \begin{columns}[T,onlytextwidth]
    \column{0.54\textwidth}
    \begin{block}{shell 中重新启动 application}
\begin{lstlisting}[
  basicstyle=\ttfamily\fontsize{5}{6}\selectfont\color{terminalfg},
  backgroundcolor=\color{terminalbg},
  frame=single,
  rulecolor=\color{terminalborder},
  framesep=1.8mm,
  keepspaces=true,
  columns=fullflexible,
  escapeinside={(*@}{@*)},
  language={}
]
(app_demo@ERICKSUN-MC1)3> (*@\textcolor{primary}{application:which\_applications().}@*)
[{(*@\textcolor{erlblue}{observer\_cli}@*), "...", "1.8.8"},
 {(*@\textcolor{erlblue}{recon}@*), "Diagnostic tools ...", "2.5.6"},
 {inets, "INETS  CXC 138 49", "9.1.0.4"},
 {ssl, "Erlang/OTP SSL application", "11.1.4.10"},
 {crypto, "CRYPTO", "5.4.2.3"},
 {stdlib, "ERTS  CXC 138 10", "5.2.3.5"},
 {(*@\textcolor{primary}{kernel}@*), "ERTS  CXC 138 10", "9.2.4.10"}]

(app_demo@ERICKSUN-MC1)4> (*@\textcolor{primary}{application:start(full\_otp\_app).}@*)
(*@\textcolor{accent}{[full\_otp\_app] starting...}@*)
(*@\textcolor{accent}{[kv\_server] init: pid=<0.317.0>}@*)
(*@\textcolor{accent}{[event\_bus\_handler] init}@*)
(*@\textcolor{accent}{[event\_bus] started with default handler}@*)
(*@\textcolor{accent}{[traffic\_light] init: starting at red}@*)
(*@\textcolor{accent}{[data\_store] init: pid=<0.321.0>}@*)
(*@\textcolor{accent}{[data\_cache] init: ETS cache table created}@*)
(*@\textcolor{accent}{ok}@*)
\end{lstlisting}
    \end{block}
    \column{0.40\textwidth}
    \begin{block}{截图解析}
      \begin{itemize}
        \item \code{which\_applications()} 中已经没有 \code{full\_otp\_app} — stop 确实清除了它。
        \item \code{application:start(full\_otp\_app).} 触发 callback module，再次启动 top supervisor 和 workers。
        \item 日志再次出现：\code{kv\_server}、\code{event\_bus}、\code{traffic\_light}、\code{data\_store}、\code{data\_cache} 又被初始化。
      \end{itemize}
    \end{block}
    \vspace{0.15cm}
    \centering
    \good{application:stop 之后，application:start 可以重新建立完整的进程树。}
  \end{columns}
\end{frame}

\begin{frame}[fragile]{第十张图：application 模块的入口很集中}
  \begin{columns}[T,onlytextwidth]
    \column{0.54\textwidth}
    \begin{block}{shell 中输入 application: + Tab}
\begin{lstlisting}[
  basicstyle=\ttfamily\scriptsize\color{terminalfg},
  backgroundcolor=\color{terminalbg},
  frame=single,
  rulecolor=\color{terminalborder},
  framesep=2.2mm,
  keepspaces=true,
  columns=fullflexible,
  escapeinside={(*@}{@*)},
  language={}
]
(app_demo@ERICKSUN-MC1)6> (*@\textcolor{primary}{application:}@*)
behaviour_info/1     ensure_all_started/1,2
ensure_started/1     get_all_env/1
get_all_key/1        get_application/1,2
get_env/1,2,3        get_key/2,3
get_supervisor/1     info/1
load/1               loaded_applications/0
permit/2             set_env/3,4
(*@\textcolor{accent}{start/1,2}@*)          start_boot/1,2
start_type/0         (*@\textcolor{accent}{stop/1}@*)
takeover/2           unload/1
unset_env/2,3        (*@\textcolor{accent}{which\_applications/0}@*)
\end{lstlisting}
    \end{block}
    \column{0.42\textwidth}
    \begin{exampleblock}{入门阶段最常用}
      \begin{itemize}
        \item \code{start/1}
        \item \code{stop/1}
        \item \code{ensure\_all\_started/1}
        \item \code{which\_applications/0}
        \item \code{get\_env/2} / \code{set\_env/3}
      \end{itemize}
    \end{exampleblock}
    \vspace{0.2cm}
    \begin{alertblock}{为什么说使用简单}
      用户入口集中在配置文件和少量 API；复杂实现被 OTP 和工具链包装了。
    \end{alertblock}
  \end{columns}
\end{frame}

% ── Summary ────────────────────────────────────────────────────────
\begin{frame}{runtime application 不是 main 函数}
  \begin{columns}[T,onlytextwidth]
    \column{0.48\textwidth}
    \begin{block}{main 函数思维}
      \begin{itemize}
        \item 从入口执行一段逻辑
        \item 执行完可能就退出
        \item 重点是控制流
      \end{itemize}
    \end{block}
    \column{0.48\textwidth}
    \begin{exampleblock}{OTP application 思维}
      \begin{itemize}
        \item 从入口建立运行时结构
        \item 启动 top supervisor
        \item supervisor 管理长期 process
        \item 重点是生命周期和容错结构
      \end{itemize}
    \end{exampleblock}
  \end{columns}
  \vspace{0.25cm}
  \centering
  \good{runtime application 是启动边界，真正长期工作的是它启动出来的 supervisor 和 worker。}
\end{frame}

\begin{frame}{本集小结}
  \begin{columns}[T,onlytextwidth]
    \column{0.50\textwidth}
    \begin{block}{runtime application 的特点}
      \begin{itemize}
        \item 有 application resource file
        \item 通常有 \code{mod} 字段
        \item 通常有 application callback module
        \item \code{start/2} 通常启动 top supervisor
        \item 启动后拥有 supervision tree
      \end{itemize}
    \end{block}
    \column{0.46\textwidth}
    \begin{exampleblock}{这一集用截图证明了什么}
      \begin{itemize}
        \item rebar3 认识 application
        \item \code{.app.src} 声明启动入口和依赖
        \item shell API 能 start / stop application
        \item Observer 能看到 app 的真实进程树
      \end{itemize}
    \end{exampleblock}
  \end{columns}
  \vspace{0.25cm}
  \centering
  \warn{下一集：library application。它也属于 OTP application，但通常没有自己的独立进程树。}
\end{frame}

\end{document}
